\begingroup\small
Let a triangle have sides $a,b,c$, semiperimeter $s$, area $K$, inradius
$r$, circumradius $R$, and median $m_a$.
\begin{align*}
 s&=\frac{a+b+c}{2}, &
 K&=\sqrt{s(s-a)(s-b)(s-c)}=rs=\frac{abc}{4R},\\
 m_a&=\frac12\sqrt{2b^2+2c^2-a^2}, &
 h_a&=\frac{2K}{a},\\
 \ell_a&=\frac{\sqrt{bc((b+c)^2-a^2)}}{b+c}, &
 OI^2&=R(R-2r).
\end{align*}
Centroid divides each median $2:1$. If the angle bisector from $A$ meets
$BC$ at $D$, then $BD:DC=AB:AC$. For a cyclic quadrilateral,
\[
 AC\cdot BD=AB\cdot CD+BC\cdot DA,\qquad
 K=\sqrt{(s-a)(s-b)(s-c)(s-d)}.
\]
For vectors and a nondegenerate polygon $p_0,\ldots,p_{n-1}$ with
$p_n=p_0$, let $S=\sum_i p_i\times p_{i+1}$ (signed doubled area). Then
\begin{align*}
 d(p,\overleftrightarrow{ab})
   &=\frac{|(b-a)\times(p-a)|}{|b-a|},\\
 \operatorname{proj}_{v}(u)&=v\frac{u\cdot v}{v\cdot v},\\
 K&=\frac{|S|}{2},\\
 G&=\frac1{3S}\sum_i(p_i+p_{i+1})(p_i\times p_{i+1}).
\end{align*}
For a circle of radius $r$ and central angle $\theta$ (radians),
\begin{gather*}
 L=r\theta,\qquad A_{\rm sector}=\frac{r^2\theta}{2},\qquad
 A_{\rm segment}=\frac{r^2}{2}(\theta-\sin\theta),\\
 \operatorname{chord}=2r\sin\frac\theta2,\qquad
 d_{\rm center\text{-}chord}=r\cos\frac\theta2.
\end{gather*}
For a regular $n$-gon of side $a$,
\[
 R=\frac{a}{2\sin(\pi/n)},\quad r=\frac{a}{2\tan(\pi/n)},\quad
 K=\frac{na^2}{4\tan(\pi/n)}.
\]
For lattice points $A,B$, a closed segment contains
$\gcd(|\Delta x|,|\Delta y|)+1$ lattice points; hence a lattice polygon has
$B=\sum_i\gcd(|\Delta x_i|,|\Delta y_i|)$ boundary points. Pick's theorem:
$K=I+B/2-1$. Common volumes:
\begin{align*}
 V_{\rm sphere}&=\frac43\pi r^3, & V_{\rm cone}&=\frac13\pi r^2h,\\
 V_{\rm cylinder}&=\pi r^2h, & V_{\rm pyramid}&=\frac13Bh.
\end{align*}
For integer coordinates, use a wide type for cross products. For floating
geometry, compare $|x|\le\varepsilon$ and avoid mixing incompatible epsilons.
\endgroup
